%% ------------------------------------------------------------------
%% Problema 5
%% ------------------------------------------------------------------
\section{Problema 5. M\'aximo de una matriz}

\subsection*{Enunciado}
Dado un arreglo bidimensional $A$ de tama\~no $N \times M$ ($N$ filas y $M$
columnas), se requiere dise\~nar un algoritmo concurrente que determine el valor
m\'aximo global de la matriz utilizando paralelismo de datos, modelado con
grafos de dependencia, condiciones de Bernstein y primitivas
\texttt{COBEGIN--COEND}.

\subsection*{An\'alisis del algoritmo}
La matriz se particiona horizontalmente asignando una fila completa a cada
proceso concurrente:
\begin{enumerate}[leftmargin=1.4cm]
    \item \textbf{Fase Concurrente ($S_0, S_1, \dots, S_{N-1}$):} Cada proceso
    $S_i$ calcula el valor m\'aximo de su fila correspondiente y lo almacena en
    la posici\'on $i$ de un vector auxiliar unidimensional $\text{max\_fila}$:
    \[
        \text{max\_fila}[i] = \max_{0 \le j < M} \{ A[i][j] \}.
    \]
    \item \textbf{Fase de Reducci\'on Secuencial ($S_{\text{final}}$):} Un proceso
    final recorre el arreglo auxiliar $\text{max\_fila}$ de longitud $N$ y
    extrae el m\'aximo global:
    \[
        \text{max\_global} = \max_{0 \le i < N} \{ \text{max\_fila}[i] \}.
    \]
\end{enumerate}

\subsection*{Condiciones de Bernstein}
Para cada proceso se identifican los conjuntos de variables que lee ($L$) y que
escribe ($E$):

\begin{center}
\begin{tabular}{lll}
\toprule
Proceso / Tarea & Lee $L$ & Escribe $E$ \\
\midrule
$S_i$ (fila $i$, con $0 \le i < N$) & $\{ A[i][j] \mid 0 \le j < M \}$ & $\{ \text{max\_fila}[i] \}$ \\
$S_{\text{final}}$ (reducci\'on)     & $\{ \text{max\_fila}[i] \mid 0 \le i < N \}$ & $\{ \text{max\_global} \}$ \\
\bottomrule
\end{tabular}
\end{center}

Verificaci\'on de independencia:
\begin{itemize}[leftmargin=1.4cm]
    \item \textbf{Entre filas independientes ($S_a \parallel S_b$, con $a \neq b$):}
    \[
        L(S_a) \cap E(S_b) = \emptyset, \qquad
        E(S_a) \cap L(S_b) = \emptyset, \qquad
        E(S_a) \cap E(S_b) = \emptyset.
    \]
    Cada proceso lee elementos de su propia fila y escribe en una posici\'on
    desacoplada del arreglo. Al cumplirse simult\'aneamente las tres
    condiciones, todos los procesos $\{S_0, S_1, \dots, S_{N-1}\}$ pueden
    ejecutarse en paralelo.
    \item \textbf{Dependencia con el proceso final ($S_i \rightarrow S_{\text{final}}$):}
    \[
        E(S_i) \cap L(S_{\text{final}}) = \{ \text{max\_fila}[i] \} \neq \emptyset.
    \]
    Existe una dependencia estricta de tipo RAW (\emph{Read-After-Write}). Por
    ende, $S_{\text{final}}$ no puede ejecutarse concurrentemente con las filas
    y requiere una barrera de sincronizaci\'on previa.
\end{itemize}

\subsection*{Grafo de dependencias}
Los procesos de fila se bifurcan en paralelo y convergen mediante una barrera
de sincronizaci\'on (\texttt{COEND}) antes de ejecutar la reducci\'on secuencial:

\begin{center}
\begin{tikzpicture}[node distance=1.4cm]
    \node[task] (inicio) {Inicio};
    \node[wide] (s0) at (-5.4,-2.0) {S0: max\_fila[0]};
    \node[wide] (s1) at (-1.8,-2.0) {S1: max\_fila[1]};
    \node[font=\large] (dots) at (1.8,-2.0) {$\cdots$};
    \node[wide] (sn) at (5.4,-2.0) {S\_N-1: max\_fila[N-1]};

    \node[wide] (sfinal) at (0,-3.8) {S\_final: max\_global};
    \node[task] (fin) at (0,-5.2) {Fin};

    \node[font=\scriptsize\itshape, text=red!70!black, align=center,
          anchor=west] at (1.9,-3.8)
          {Barrera de sincronizaci\'on\\(\texttt{COEND})};

    \draw[dep] (inicio) -- (s0);
    \draw[dep] (inicio) -- (s1);
    \draw[dep] (inicio) -- (sn);

    \draw[dep] (s0) -- (sfinal);
    \draw[dep] (s1) -- (sfinal);
    \draw[dep] (sn) -- (sfinal);

    \draw[dep] (sfinal) -- (fin);
\end{tikzpicture}
\end{center}

El camino cr\'itico es $\text{Inicio} \rightarrow S_i \rightarrow S_{\text{final}} \rightarrow \text{Fin}$ y el paralelismo m\'aximo es $N$.

\subsection*{Soluci\'on con COBEGIN--COEND}
\begin{lstlisting}
// Entrada: A[N][M] (Matriz de NxM)
// Auxiliar: max_fila[N]
// Salida: max_global

COBEGIN
    S_0;     // max_fila[0]   = max(A[0][0..M-1])
    S_1;     // max_fila[1]   = max(A[1][0..M-1])
    ...
    S_N_1;   // max_fila[N-1] = max(A[N-1][0..M-1])
COEND

S_final;     // max_global    = max(max_fila[0..N-1])
escribir max_global
\end{lstlisting}

\noindent
La b\'usqueda del m\'aximo en cada fila es totalmente independiente, por lo que
se ejecuta dentro del bloque \texttt{COBEGIN}. La variable \texttt{max\_global}
representa la consolidaci\'on de los resultados intermedios y s\'olo puede
calcularse una vez que la barrera \texttt{COEND} garantiza que todas las filas
concluyeron su trabajo.
